% =====================================================================
%  Reporte: Pipeline CI/CD para API REST con Docker, GitHub Actions y AWS EC2
%  Compilar con pdfLaTeX (Overleaf: Menu -> Compiler -> pdfLaTeX)
%  Las capturas van en la carpeta figuras/ con el nombre que indica cada recuadro.
% =====================================================================
\documentclass[12pt,letterpaper]{article}

% ---------- Idioma y codificación ----------
\usepackage[utf8]{inputenc}
\usepackage[T1]{fontenc}
\usepackage[spanish,es-tabla,es-noshorthands]{babel}
\usepackage{lmodern}
\usepackage{microtype}

% ---------- Formato de página ----------
\usepackage[margin=2.5cm]{geometry}
\usepackage{setspace}
\onehalfspacing
\setlength{\parskip}{6pt}
\usepackage{fancyhdr}
\pagestyle{fancy}
\fancyhf{}
\fancyhead[L]{\small Pipeline CI/CD para API REST}
\fancyhead[R]{\small UTEQ}
\fancyfoot[C]{\thepage}

% ---------- Figuras, tablas y código ----------
\usepackage{graphicx}
\usepackage{float}
\usepackage[labelfont=bf]{caption}
\usepackage{booktabs}
\usepackage{tabularx}
\usepackage{xcolor}
\usepackage{listings}
\usepackage{tikz}
\usetikzlibrary{positioning,arrows.meta}
\usepackage[hidelinks]{hyperref}

\definecolor{fondo}{RGB}{246,248,250}
\definecolor{comentario}{RGB}{106,115,125}
\definecolor{clave}{RGB}{0,92,197}
\definecolor{cadena}{RGB}{3,47,98}

\lstset{
  basicstyle=\ttfamily\footnotesize,
  backgroundcolor=\color{fondo},
  commentstyle=\color{comentario}\itshape,
  keywordstyle=\color{clave}\bfseries,
  stringstyle=\color{cadena},
  frame=single, rulecolor=\color{black!20},
  breaklines=true, columns=fullflexible,
  numbers=left, numberstyle=\tiny\color{comentario},
  showstringspaces=false,
  inputencoding=utf8, extendedchars=true,
  literate={á}{{\'a}}1 {é}{{\'e}}1 {í}{{\'i}}1 {ó}{{\'o}}1 {ú}{{\'u}}1
           {Á}{{\'A}}1 {É}{{\'E}}1 {Í}{{\'I}}1 {Ó}{{\'O}}1 {Ú}{{\'U}}1
           {ñ}{{\~n}}1 {Ñ}{{\~N}}1 {¿}{{?`}}1 {¡}{{!`}}1
}
\renewcommand{\lstlistingname}{Código}
\renewcommand{\lstlistlistingname}{Índice de códigos}

% Inserta una captura; si aún no existe muestra un recuadro con el nombre esperado.
% Uso: \captura{archivo.png}{Pie de figura}{etiqueta}
\newcommand{\captura}[3]{%
  \begin{figure}[H]
    \centering
    \IfFileExists{figuras/#1}{%
      \includegraphics[width=0.92\textwidth,height=0.42\textheight,keepaspectratio]{figuras/#1}%
    }{%
      \fbox{\parbox[c][5cm][c]{0.85\textwidth}{\centering\small
        Captura pendiente\\[4pt]\texttt{figuras/#1}}}%
    }
    \caption{#2}
    \label{#3}
  \end{figure}}

% ---------- Datos de la portada (EDITAR) ----------
\newcommand{\universidad}{Universidad Tecnológica de Querétaro}
\newcommand{\carrera}{[Nombre de la carrera]}
\newcommand{\materia}{[Nombre de la asignatura]}
\newcommand{\titulo}{Pipeline CI/CD automatizado para una API REST con Docker, GitHub Actions y AWS EC2}
\newcommand{\profesor}{[Nombre del profesor]}
\newcommand{\grupo}{[Grupo]}
\newcommand{\fecha}{[Fecha de entrega]}
\newcommand{\repositorio}{https://github.com/[usuario]/mi-webapp}
\newcommand{\ipec}{[IP\_EC2]}

% Título de la bibliografía
\addto\captionsspanish{\renewcommand{\refname}{Fuentes de información}}

\begin{document}

% =====================================================================
%  PORTADA
% =====================================================================
\begin{titlepage}
  \thispagestyle{empty}
  \centering
  \IfFileExists{figuras/logo-uteq.png}{%
    \includegraphics[width=0.35\textwidth]{figuras/logo-uteq.png}\par
  }{%
    \fbox{\parbox[c][3cm][c]{0.35\textwidth}{\centering\small Logo institucional\\\texttt{figuras/logo-uteq.png}}}\par
  }
  \vspace{1cm}
  {\LARGE\bfseries \universidad\par}
  \vspace{0.4cm}
  {\large \carrera\par}
  \vspace{1.5cm}
  {\large \materia\par}
  \vspace{1cm}
  \rule{\textwidth}{1pt}\par\vspace{0.4cm}
  {\Large\bfseries \titulo\par}
  \vspace{0.4cm}\rule{\textwidth}{1pt}\par
  \vspace{1.5cm}
  {\large\bfseries Integrantes\par}
  \vspace{0.3cm}
  \begin{tabular}{ll}
    {[Nombre completo]} & 2023371111 \\
    {[Nombre completo]} & {[Matrícula]} \\
    {[Nombre completo]} & {[Matrícula]} \\
  \end{tabular}\par
  \vfill
  {\large \textbf{Profesor:} \profesor\par}
  \vspace{0.3cm}
  {\large \textbf{Grupo:} \grupo\par}
  \vspace{0.8cm}
  {\large Querétaro, Qro., \fecha\par}
\end{titlepage}

% =====================================================================
%  ÍNDICES
% =====================================================================
\tableofcontents
\newpage
\listoffigures
\listoftables
\lstlistoflistings
\newpage

% =====================================================================
%  INTRODUCCIÓN (mínimo una página completa)
% =====================================================================
\section{Introducción}

En el desarrollo de software actual, escribir código que funcione ya no es suficiente: las organizaciones necesitan entregar nuevas versiones de sus aplicaciones de manera frecuente, confiable y sin interrumpir el servicio a los usuarios. Durante muchos años, el proceso de llevar un cambio desde la computadora del programador hasta el servidor de producción fue una tarea manual que implicaba compilar, ejecutar pruebas, copiar archivos y reiniciar servicios. Cada uno de estos pasos manuales representaba una oportunidad de error humano, provocaba diferencias entre el ambiente de desarrollo y el de producción, y hacía que las entregas fueran lentas y riesgosas.

Como respuesta a esta problemática surgió la cultura DevOps, que busca integrar el trabajo de los equipos de desarrollo y de operaciones mediante la automatización y la retroalimentación continua. Dos de sus prácticas centrales son la Integración Continua (CI, por sus siglas en inglés), en la que cada cambio que se sube al repositorio se compila y se prueba automáticamente \cite{fowler}, y la Entrega o Despliegue Continuo (CD), en la que los cambios que superan las pruebas se empaquetan y se publican en el servidor sin intervención manual \cite{humble}. Al conjunto de etapas automatizadas que recorre el código se le conoce como \emph{pipeline} CI/CD.

La contenedorización es otra pieza clave de este enfoque. Docker permite empaquetar una aplicación junto con su entorno de ejecución, sus bibliotecas y su configuración en una imagen inmutable que se comporta igual en la computadora del desarrollador, en el servidor de integración y en la nube \cite{merkel,dockerdocs}. Estas imágenes se almacenan en un registro, como Docker Hub, desde donde cualquier servidor autorizado puede descargarlas. Por su parte, los proveedores de nube como Amazon Web Services ofrecen servidores virtuales bajo demanda, como las instancias EC2, que pueden configurarse en minutos y pagarse solo por el tiempo de uso \cite{awsec2}.

En este proyecto integrador se desarrolló una API REST en Node.js con el framework Express y una base de datos SQLite, siguiendo el estilo arquitectónico REST propuesto por Fielding \cite{fielding}. La API expone seis endpoints: uno de verificación de salud y cinco para la gestión de usuarios (listar, consultar, crear, actualizar y eliminar). Sobre esta aplicación se construyó un pipeline en GitHub Actions \cite{githubactions} que, ante cada \emph{push} o \emph{pull request} a la rama principal, ejecuta las pruebas automatizadas con Jest y Supertest, verifica que la cobertura de código sea al menos del 70\%, construye la imagen Docker, la publica en Docker Hub con las etiquetas \texttt{latest} y el identificador del commit, y finalmente se conecta por SSH a una instancia EC2 con Ubuntu Server para reemplazar el contenedor anterior por la nueva versión.

Un aspecto fundamental del proyecto es la seguridad: ninguna contraseña, token, dirección IP o llave privada se encuentra en el código fuente. Todos los datos sensibles se almacenan como \emph{GitHub Secrets}, que el pipeline lee en tiempo de ejecución y que GitHub oculta automáticamente en los registros. Asimismo, el contenedor se ejecuta con un usuario sin privilegios y el pipeline verifica el estado de la API después de cada despliegue, restaurando la versión anterior en caso de falla.

El presente documento describe la arquitectura de la solución, el proceso de configuración de cada componente y los resultados obtenidos, los cuales se evidencian mediante capturas de pantalla numeradas. Finalmente, se presenta una demostración en la que un cambio de una sola línea en el código se refleja automáticamente en el servidor público, lo que comprueba el funcionamiento completo del ciclo de integración y despliegue continuo.

\subsection{Objetivo general}
Implementar un flujo de integración y despliegue continuo que automatice las pruebas, la generación de imágenes de contenedor y la publicación en AWS EC2 de una API REST.

\subsection{Objetivos específicos}
\begin{itemize}
  \item Desarrollar una API REST con seis endpoints funcionales y validación de datos.
  \item Implementar pruebas de integración con una cobertura de código mínima del 70\%.
  \item Contenedorizar la aplicación mediante un Dockerfile optimizado y un archivo \texttt{.dockerignore}.
  \item Configurar un workflow de GitHub Actions que pruebe, construya y publique la imagen en Docker Hub.
  \item Desplegar automáticamente la aplicación en una instancia EC2 por medio de SSH.
  \item Gestionar todos los datos sensibles a través de GitHub Secrets.
\end{itemize}

% =====================================================================
%  MARCO TEÓRICO
% =====================================================================
\section{Marco teórico}

\paragraph{API REST.} Interfaz de programación que sigue el estilo arquitectónico REST, en el que los recursos se identifican mediante URL y se manipulan con los métodos del protocolo HTTP (GET, POST, PUT y DELETE), intercambiando representaciones en formato JSON \cite{fielding}.

\paragraph{Integración Continua (CI).} Práctica en la que los integrantes de un equipo integran su trabajo con frecuencia y cada integración se verifica con una compilación y pruebas automáticas, con el fin de detectar errores lo antes posible \cite{fowler}.

\paragraph{Despliegue Continuo (CD).} Extensión de la integración continua en la que todo cambio que supera las pruebas se libera automáticamente al ambiente de producción \cite{humble}.

\paragraph{Docker.} Plataforma de contenedores que empaqueta una aplicación y sus dependencias en una imagen portable. Una imagen se define en un archivo \texttt{Dockerfile} y se ejecuta como un contenedor aislado \cite{merkel,dockerdocs}.

\paragraph{GitHub Actions.} Servicio de automatización integrado en GitHub que ejecuta \emph{workflows} definidos en archivos YAML como respuesta a eventos del repositorio, como un \emph{push} o un \emph{pull request} \cite{githubactions}.

\paragraph{Amazon EC2.} Servicio de Amazon Web Services que proporciona servidores virtuales en la nube. Su acceso de red se controla mediante \emph{Security Groups}, que funcionan como un firewall a nivel de instancia \cite{awsec2}.

\paragraph{Cobertura de código.} Métrica que indica el porcentaje de sentencias, ramas, funciones y líneas del código que se ejecutan durante las pruebas automatizadas \cite{jest}.

% =====================================================================
%  DESARROLLO
% =====================================================================
\section{Desarrollo}

\subsection{Arquitectura de la solución}

La Figura~\ref{fig:arquitectura} muestra el flujo completo. El desarrollador sube un cambio a la rama \texttt{main}; GitHub Actions ejecuta tres trabajos en secuencia (\emph{test}, \emph{build} y \emph{deploy}) y, si alguno falla, los siguientes no se ejecutan. En un \emph{pull request} solo se ejecutan las pruebas.

\begin{figure}[H]
  \centering
  \begin{tikzpicture}[
      node distance=0.6cm and 0.8cm,
      caja/.style={draw, rounded corners, align=center, minimum width=2.6cm, minimum height=1.1cm, font=\small, fill=blue!5},
      flecha/.style={-{Stealth}, thick}]
    \node[caja] (dev) {Desarrollador\\\texttt{git push}};
    \node[caja, right=of dev] (gh) {Repositorio\\GitHub};
    \node[caja, right=of gh] (test) {Job 1: test\\Jest + cobertura};
    \node[caja, below=of test] (build) {Job 2: build\\Docker build/push};
    \node[caja, left=of build] (hub) {Docker Hub\\\texttt{:latest} \texttt{:sha}};
    \node[caja, below=of build] (deploy) {Job 3: deploy\\SSH a EC2};
    \node[caja, left=of deploy, fill=orange!10] (ec2) {AWS EC2 Ubuntu\\Docker, puerto 80};
    \node[caja, left=of ec2, fill=green!8] (user) {Usuario\\\texttt{http://IP/api}};
    \draw[flecha] (dev) -- (gh);
    \draw[flecha] (gh) -- (test);
    \draw[flecha] (test) -- (build);
    \draw[flecha] (build) -- (hub);
    \draw[flecha] (build) -- (deploy);
    \draw[flecha] (deploy) -- (ec2);
    \draw[flecha, dashed] (hub) -- node[left, font=\scriptsize]{pull} (ec2);
    \draw[flecha] (user) -- (ec2);
  \end{tikzpicture}
  \caption{Arquitectura del pipeline CI/CD.}
  \label{fig:arquitectura}
\end{figure}

\subsection{API REST}

La API se desarrolló en Node.js 22 con Express 5 y almacena los datos en SQLite. La Tabla~\ref{tab:endpoints} resume los seis endpoints. Todas las respuestas comparten el formato \texttt{\{ statusCode, data \}}, y los errores incluyen un mensaje descriptivo.

\begin{table}[H]
  \centering
  \caption{Endpoints expuestos por la API.}
  \label{tab:endpoints}
  \begin{tabularx}{\textwidth}{c l l X}
    \toprule
    \# & Método & Ruta & Descripción \\
    \midrule
    1 & GET    & \texttt{/api/health}    & Estado de la API, mensaje y commit desplegado \\
    2 & GET    & \texttt{/api/users}     & Lista todos los usuarios \\
    3 & GET    & \texttt{/api/users/:id} & Obtiene un usuario por su identificador \\
    4 & POST   & \texttt{/api/users}     & Crea un usuario (nombre y correo) \\
    5 & PUT    & \texttt{/api/users/:id} & Actualiza un usuario existente \\
    6 & DELETE & \texttt{/api/users/:id} & Elimina un usuario \\
    \bottomrule
  \end{tabularx}
\end{table}

El Código~\ref{lst:post} muestra, como ejemplo, el endpoint de creación de usuarios: primero valida los datos recibidos y, si son correctos, los inserta en la base de datos y responde con el código 201.

\begin{lstlisting}[language=Java, caption={Endpoint POST /api/users (\texttt{src/app.js}).}, label={lst:post}]
// 4. POST /api/users - crear usuario
app.post('/api/users', async (req, res) => {
    const errors = validateUser(req.body);
    if (errors.length) return badRequest(res, errors);
    const { lastID } = await db.run(
        'INSERT INTO users (name, email) VALUES (?, ?)',
        [req.body.name, req.body.email]);
    send(res, await findUser(lastID), 201);
});
\end{lstlisting}

\subsection{Pruebas automatizadas}

Las pruebas se escribieron con Jest y Supertest en el archivo \texttt{tests/api.test.js}. Se ejecutan sobre una base de datos en memoria y comprueban tanto los casos exitosos de los seis endpoints como los casos de error: datos faltantes o inválidos (400), registros inexistentes (404), JSON mal formado y fallas de la base de datos (500). En el archivo \texttt{package.json} se configuró un umbral mínimo de cobertura del 70\% para sentencias, ramas, funciones y líneas; si no se alcanza, el comando \texttt{npm test} falla y el pipeline se detiene.

\subsection{Contenedorización}

El \texttt{Dockerfile} (Código~\ref{lst:docker}) utiliza una construcción en dos etapas: la primera instala las dependencias de producción y la segunda copia solo lo necesario a una imagen limpia basada en \texttt{node:22-alpine}. El contenedor se ejecuta con el usuario sin privilegios \texttt{node} e incluye un \texttt{HEALTHCHECK}. El archivo \texttt{.dockerignore} excluye \texttt{node\_modules}, archivos \texttt{.env}, llaves, registros, la base de datos local y las pruebas.

\begin{lstlisting}[caption={Dockerfile de dos etapas.}, label={lst:docker}]
FROM node:22-alpine AS deps
WORKDIR /app
RUN apk add --no-cache python3 make g++
COPY package.json package-lock.json ./
RUN npm ci --omit=dev && npm cache clean --force

FROM node:22-alpine
WORKDIR /app
ENV NODE_ENV=production PORT=3000 DB_FILE=/app/data/database.sqlite
ARG GIT_SHA=local
ENV GIT_SHA=$GIT_SHA
COPY --from=deps --chown=node:node /app/node_modules ./node_modules
COPY --chown=node:node package.json ./
COPY --chown=node:node src ./src
RUN mkdir -p /app/data && chown node:node /app/data
USER node
EXPOSE 3000
HEALTHCHECK CMD wget -qO- http://localhost:3000/api/health || exit 1
CMD ["node", "src/server.js"]
\end{lstlisting}

\subsection{Workflow de GitHub Actions}

El archivo \texttt{.github/workflows/main.yml} define los tres trabajos del pipeline. El Código~\ref{lst:workflow} muestra un extracto con los disparadores y la publicación de la imagen con las etiquetas \texttt{latest} y el hash del commit.

\begin{lstlisting}[caption={Extracto de \texttt{.github/workflows/main.yml}.}, label={lst:workflow}]
on:
  push:
    branches: [main]
  pull_request:
    branches: [main]

jobs:
  test:
    steps:
      - run: npm ci
      - run: npm test            # pruebas + cobertura (minimo 70%)
  build:
    needs: test
    steps:
      - uses: docker/login-action@v3
        with:
          username: ${{ secrets.DOCKERHUB_USERNAME }}
          password: ${{ secrets.DOCKERHUB_TOKEN }}
      - uses: docker/build-push-action@v6
        with:
          push: true
          tags: |
            ${{ env.IMAGE }}:latest
            ${{ env.IMAGE }}:${{ github.sha }}
  deploy:
    needs: build
    steps:
      - uses: appleboy/ssh-action@v1.2.0
        with:
          host: ${{ secrets.EC2_HOST }}
          username: ${{ secrets.EC2_USER }}
          key: ${{ secrets.EC2_SSH_KEY }}
\end{lstlisting}

En el servidor, el trabajo de despliegue ejecuta el script del Código~\ref{lst:deploy}: descarga la nueva imagen antes de detener la anterior para reducir el tiempo sin servicio, levanta el contenedor en el puerto 80 y verifica el endpoint de salud. Si la nueva versión no responde, restaura la anterior.

\begin{lstlisting}[language=bash, caption={Script de despliegue ejecutado en la instancia EC2.}, label={lst:deploy}]
docker pull "$IMAGE:$SHA"
docker stop api || true
docker rm api || true
docker run -d --name api --restart unless-stopped \
  -p 80:3000 -v api-data:/app/data "$IMAGE:$SHA"
curl -fsS http://localhost/api/health
\end{lstlisting}

\subsection{Gestión de datos sensibles}

La Tabla~\ref{tab:secrets} enumera los secretos configurados en el repositorio. Ninguno de estos valores aparece en el código fuente, y GitHub los oculta automáticamente en los registros del pipeline.

\begin{table}[H]
  \centering
  \caption{GitHub Secrets utilizados por el pipeline.}
  \label{tab:secrets}
  \begin{tabularx}{\textwidth}{l X}
    \toprule
    Secreto & Contenido \\
    \midrule
    \texttt{DOCKERHUB\_USERNAME} & Usuario de Docker Hub \\
    \texttt{DOCKERHUB\_TOKEN}    & Personal Access Token de Docker Hub con permiso de lectura y escritura \\
    \texttt{EC2\_HOST}           & Dirección IP pública de la instancia EC2 \\
    \texttt{EC2\_USER}           & Usuario del sistema operativo (\texttt{ubuntu}) \\
    \texttt{EC2\_SSH\_KEY}       & Contenido de la llave privada \texttt{.pem} \\
    \bottomrule
  \end{tabularx}
\end{table}

% =====================================================================
%  RESULTADOS
% =====================================================================
\section{Resultados}

En esta sección se presentan las evidencias del funcionamiento de cada etapa, en el mismo orden en que se configuró el proyecto.

\subsection{Pruebas y contenedor en el ambiente local}

Antes de automatizar el proceso se verificó que la aplicación funcionara localmente. En la Figura~\ref{fig:pruebas} se observa la ejecución del comando \texttt{npm test}: las 14 pruebas se completaron correctamente y el reporte de cobertura alcanzó 95.19\% de sentencias, 86.53\% de ramas, 96.77\% de funciones y 96.20\% de líneas, valores superiores al umbral exigido del 70\%.

\captura{fig01-pruebas-locales.png}{Ejecución local de las pruebas y reporte de cobertura.}{fig:pruebas}

Posteriormente se construyó la imagen Docker y se ejecutó un contenedor local para comprobar que la API respondiera dentro de él, como se muestra en la Figura~\ref{fig:dockerlocal}.

\captura{fig02-docker-local.png}{Construcción y ejecución local de la imagen Docker.}{fig:dockerlocal}

\subsection{Repositorio y secretos}

El código fuente se publicó en un repositorio de GitHub (Figura~\ref{fig:repo}) que contiene la API, las pruebas, el \texttt{Dockerfile}, el \texttt{.dockerignore}, el workflow y la documentación en \texttt{README.md}. En la Figura~\ref{fig:secrets} se aprecian los cinco secretos configurados; GitHub muestra únicamente sus nombres y nunca su valor.

\captura{fig03-repositorio.png}{Repositorio del proyecto en GitHub.}{fig:repo}

\captura{fig04-secrets.png}{Secretos configurados en GitHub Actions.}{fig:secrets}

\subsection{Infraestructura en AWS EC2}

Se creó una instancia EC2 con Ubuntu Server (Figura~\ref{fig:ec2}). Su Security Group permite el tráfico entrante en el puerto 22 para la conexión SSH y en el puerto 80 para HTTP (Figura~\ref{fig:sg}). Después se instaló Docker mediante el script \texttt{scripts/setup-ec2.sh}, como se observa en la Figura~\ref{fig:ec2docker}.

\captura{fig05-ec2-instancia.png}{Instancia EC2 con Ubuntu Server en ejecución.}{fig:ec2}

\captura{fig06-security-group.png}{Reglas de entrada del Security Group (puertos 22 y 80).}{fig:sg}

\captura{fig07-ec2-docker.png}{Docker instalado en la instancia EC2.}{fig:ec2docker}

\subsection{Ejecución del pipeline}

Al realizar un \emph{push} a la rama \texttt{main}, GitHub Actions ejecutó los tres trabajos de forma secuencial y todos finalizaron con éxito, como se muestra en la Figura~\ref{fig:pipeline}. En el registro del trabajo de pruebas se muestra el resumen de cobertura (Figura~\ref{fig:cobpipeline}), lo que confirma que la validación de calidad forma parte del proceso automático.

\captura{fig08-pipeline.png}{Ejecución exitosa del pipeline con sus tres trabajos.}{fig:pipeline}

\captura{fig09-cobertura-pipeline.png}{Resumen de cobertura en los registros del pipeline.}{fig:cobpipeline}

El trabajo \emph{build} publicó la imagen en Docker Hub con dos etiquetas: \texttt{latest} y el hash completo del commit (Figura~\ref{fig:dockerhub}). Esta última permite identificar con exactitud qué versión del código contiene cada imagen y regresar a una versión anterior si fuera necesario.

\captura{fig10-dockerhub.png}{Etiquetas de la imagen publicadas en Docker Hub.}{fig:dockerhub}

Finalmente, el trabajo \emph{deploy} se conectó a la instancia por SSH, descargó la imagen, reemplazó el contenedor y verificó el endpoint de salud (Figura~\ref{fig:deploy}).

\captura{fig11-deploy-log.png}{Registro del despliegue en la instancia EC2.}{fig:deploy}

\subsection{API pública}

Una vez desplegada, la API respondió desde la dirección pública de la instancia, \texttt{http://\ipec/api/health} (Figura~\ref{fig:publica}). En la Figura~\ref{fig:crud} se observan las operaciones de creación, consulta, actualización y eliminación de usuarios realizadas contra el servidor en la nube.

\captura{fig12-api-publica.png}{Respuesta del endpoint de salud desde la IP pública de EC2.}{fig:publica}

\captura{fig13-crud.png}{Operaciones CRUD sobre \texttt{/api/users} en el servidor público.}{fig:crud}

\subsection{Demostración del despliegue automático}

Para comprobar el ciclo completo se modificó una sola línea del archivo \texttt{src/config.js}, correspondiente al mensaje que devuelve la API, y se realizó un \emph{push} (Figura~\ref{fig:cambio}). Sin ninguna intervención manual, el pipeline ejecutó las pruebas, publicó una nueva imagen y actualizó el servidor. En la Figura~\ref{fig:demo} se aprecia que el endpoint \texttt{/api/health} devuelve el nuevo mensaje y el identificador del nuevo commit.

\captura{fig14-demo-cambio.png}{Cambio de una línea del código y \emph{push} a la rama main.}{fig:cambio}

\captura{fig15-demo-resultado.png}{La API pública muestra el nuevo mensaje tras el despliegue automático.}{fig:demo}

% =====================================================================
%  CONCLUSIONES (mínimo dos párrafos)
% =====================================================================
\section{Conclusiones}

La implementación de este proyecto permitió comprobar que la automatización del ciclo de vida del software reduce de forma importante el esfuerzo y el riesgo asociados a la publicación de nuevas versiones. Una vez configurado el pipeline, basta un \emph{push} a la rama principal para que el código se pruebe, se empaquete y se despliegue en la nube en cuestión de minutos, sin pasos manuales. Las pruebas automatizadas y el umbral de cobertura del 70\% funcionan como un filtro de calidad: si un cambio rompe alguna funcionalidad, el pipeline se detiene antes de construir la imagen y el servidor de producción no se ve afectado.

El uso de Docker resolvió el problema de las diferencias entre ambientes, ya que la misma imagen que se probó y se publicó en Docker Hub es la que se ejecuta en la instancia EC2. Etiquetar cada imagen con el hash del commit aporta trazabilidad, pues permite saber exactamente qué versión del código está en producción y regresar a una anterior si es necesario. Del mismo modo, la verificación del endpoint de salud después de cada despliegue y la restauración automática de la versión previa aumentan la confiabilidad del servicio.

En materia de seguridad, el proyecto mostró la importancia de separar la configuración sensible del código fuente. Gracias a GitHub Secrets, el repositorio puede ser público sin exponer credenciales, direcciones IP ni llaves privadas. Como trabajo futuro se propone incorporar un balanceador o proxy inverso para lograr despliegues sin ningún segundo de inactividad, utilizar una base de datos administrada para separar los datos del servidor de aplicación, y restringir el acceso SSH mediante mecanismos como AWS Systems Manager.

% =====================================================================
%  FUENTES DE INFORMACIÓN
% =====================================================================
\begin{thebibliography}{9}

\bibitem{fowler}
Fowler, M. (2006). \emph{Continuous Integration}. martinfowler.com. Recuperado el [fecha] de \url{https://martinfowler.com/articles/continuousIntegration.html}

\bibitem{humble}
Humble, J. y Farley, D. (2010). \emph{Continuous Delivery: Reliable Software Releases through Build, Test, and Deployment Automation}. Addison-Wesley.

\bibitem{fielding}
Fielding, R. T. (2000). \emph{Architectural Styles and the Design of Network-based Software Architectures} [Tesis doctoral]. University of California, Irvine.

\bibitem{merkel}
Merkel, D. (2014). Docker: Lightweight Linux containers for consistent development and deployment. \emph{Linux Journal}, 2014(239), 2.

\bibitem{dockerdocs}
Docker Inc. (s.f.). \emph{Dockerfile reference}. Recuperado el [fecha] de \url{https://docs.docker.com/reference/dockerfile/}

\bibitem{githubactions}
GitHub. (s.f.). \emph{GitHub Actions documentation}. Recuperado el [fecha] de \url{https://docs.github.com/en/actions}

\bibitem{awsec2}
Amazon Web Services. (s.f.). \emph{Amazon EC2 User Guide}. Recuperado el [fecha] de \url{https://docs.aws.amazon.com/AWSEC2/latest/UserGuide/}

\bibitem{jest}
OpenJS Foundation. (s.f.). \emph{Jest: Configuring Jest, coverageThreshold}. Recuperado el [fecha] de \url{https://jestjs.io/docs/configuration}

\end{thebibliography}

\end{document}
